\section{Conclusion}\label{sec:conclusion}

We mechanized the computational security of a hybrid post-quantum
handshake, including the ``either component'' lemmas for X-Wing, an
ML-KEM-1024 transcript-hashing combiner and composite ML-DSA/EdDSA
signatures. The machine-checked results hold. What the exercise produced
beyond them was a list of things that had been wrong: two lemmas of the
sketch, five claims of our own write-up, and one divergence between the
implementation and its specification. Each was caught by a different
mechanism: modelling the construction exactly, an independent review of the
claims, and tests aimed at what the models abstract away. Mutation tests made
it possible to say what each proof depends on.

Several directions remain open: a cofactor-aware model of X25519;
mechanizing the composition steps that we did by hand; a tight multi-user
bound for the hybrid KEM; and adaptive ephemeral and state reveal. The
specifications also disagree on the all-zero X25519 output: X-Wing specifies
no check, while RFC~9180 and RFC~10024 require one. Implementations built on
X25519 functions that reject that output will keep diverging from X-Wing
unless the specifications settle the question.

\paragraph{Acknowledgements and AI disclosure.} The CryptoVerif models,
proof scripts, mutation tests, implementation tests and a first draft of
this paper were produced with Claude (Anthropic) at the author's direction,
and the review of \cref{sec:reviewcaught} was done by a separate Claude
instance. Every machine-checked result was re-run from a clean checkout. The
author takes responsibility for the content.
